\section{Conclusion}

We introduce human and mouse benchmarks that compare two peptides from the same parent protein under the same \plainMHC{} restriction and predict which has stronger presentation evidence in a specified tissue. Peptide sequence supports this ranking in both species, and removing tissue identity substantially reduces performance. However, performance also decreases under \plainPeptideOOF{}, demonstrating that estimates from familiar tasks with cross-task reuse allowed (standard estimates) benefit from repeated peptides across tasks. The retained signal supports learnability of tissue-associated peptide ranking for previously represented tissue--antigen-presenting-molecule combinations (tissue--MHC combinations), but does not identify a specific antigen-processing mechanism. Complex model structures (architectures) do not consistently outperform strong simpler controls under peptide separation. The main contribution is therefore a biologically focused and reproducible benchmark that distinguishes tissue-associated signal from general peptide--antigen-presenting-molecule compatibility (peptide--MHC compatibility), while making the limits of unseen-peptide prediction explicit. Validation in unseen tissues, alleles, studies, and prospective samples remains necessary before biological or clinical application.
